\scanpage{223}
\mspage{139}
% archive sheet number: 222
% element: text | top, section heading | 4.6 Caracterisation des cat proto-topos
4.6. \uline{Caractérisation des cat. \unsure{proto-topos|prototopos}} (dém. \unsure{de|du} 4.3.6)

% element: text | upper third, return to earlier proof | Il est temps de revenir
\deleted{4.6.1.} Il est temps de revenir à la démonstration du th. 4.3.6, \inserted{(p. 70--71)} celle que j'avais donnée \unsure{étant|était} incomplète (et erronée --- j'affirmais que $(\mathrm{Filt}^\pi)_{\pi\geq\pi_{\unsure{1|0}}}$ est une filtration cardinale, ce qui est faux si on n'admet que $\pi^{\pi_0}=\pi$ pour $\unsure{\forall|\mathrm{tt}}\pi\geq\pi_1$.

% element: text | middle upper, logical scheme lead | On avait vu l equivalence
On avait vu l'équivalence des conditions a, a$'$, b, soit PTA. Je vais donner la démonstration par le schéma logique
% element: diagram | middle, implication chain | PTA implies d implies c implies PTA
\diagram{\begin{tikzcd}[ampersand replacement=\&]
(\unsure{\alpha|a})\&\mathrm{PTA}\arrow[r,Rightarrow]\&d)\arrow[r,Rightarrow]\&c)\arrow[r,Rightarrow]\&\mathrm{PTA}.
\end{tikzcd}}
% element: text | middle, accessible proto-topos and small category | Prouvons PTA implies d
Prouvons $\mathrm{PTA}\Rightarrow d)$, i.e. que si $M$ est un proto-topos accessible, alors \deleted{\unsure{il y a|il existe}} \inserted{\unsure{existe}} \inserted{un cardinal $\pi_0$ \unsure{infini}, et} une \deleted{\unsure{petite}} \inserted{petite} catégorie $C_0$, \deleted{\unsure{qui|que}} \inserted{stable par $\unsure{\underline{\lim}|\varinjlim|\varprojlim}$ de cardinal $\leq\pi_0$} et une équivalence de catégories
% element: diagram | lower middle, equivalence 4.6.1 | M to Ind prime pi zero C zero
\diagram{\begin{tikzcd}[ampersand replacement=\&]
\text{(4.6.1)}\&M\arrow[r,"\approx"]\&\mathrm{Ind}'_{\pi_0}(C_0).
\end{tikzcd}}
% element: text | lower third, transition | Plus precisement
Plus précisément, on a ceci :

% element: text | lower third, proposition 4.6.1 and corrections | Proposition 4.6.1 Soit M
\uline{Proposition} 4.6.1. Soit $M$ un proto-topos, i.e. une \deleted{proto-topos} \inserted{$\unsure{U|\mathcal{U}|\mathscr{U}}$catégorie \unsure{et} stable par petites $\underline{\lim}$,} admettant une \deleted{petite} sous-catégorie \inserted{pleine} génératrice $S$, \inserted{\unsure{qui est une} petite catégorie} d'objets accessibles de $M$. Soit $\pi_0$ un $\unsure{*|+}$ cardinal tel que les objets de $S$ soient $\pi_0$-\unsure{accessibles|accessible}, \inserted{(i.e. $S\subset M_{\pi_0}$)} \deleted{$\unsure{C|S}$} et soit $C$ la plus petite sous-catégorie str. pleine de $M$, contenant \deleted{$\unsure{C|S}$} $S$ et stable par $\varinjlim$ de cardinal $\leq\pi_0$. Alors
% element: text | foot, first conclusion | C est equivalente
\unsure{1°)|a)} $C$ est équivalente à une petite catégorie.
